% Part: incompleteness
% Chapter: arithmetization-syntax
% Section: proofs-in-lk

\documentclass[../../../include/open-logic-section]{subfiles}

\begin{document}

\olfileid{inc}{art}{plk}
\olsection{\usetoken{P}{derivation} در~$\Log{LK}$}

\begin{explain}
برای حسابی‌سازی !!{derivation}s باید !!{derivation}s را با اعداد بازنمایی کنیم.
چون !!{derivation}s درخت‌هایی از سکوئنت‌ها هستند و هر استنتاج نیز
برچسبی دارد، بازنمایی بازگشتی آشکارترین رویکرد است: !!{derivation} را
به‌صورت چندتایی‌ای بازنمایی می‌کنیم که مؤلفه‌هایش سکوئنت پایانی،
برچسب و بازنمایی زیر!!{derivation}sی منتهی به مقدمات آخرین استنتاج‌اند.
\end{explain}

\begin{defn}
اگر $\Gamma$ دنباله‌ای متناهی از !!{sentence}s باشد،
$\Gamma=\tuple{!A_1,\dots,!A_n}$، آنگاه
$\Gn{\Gamma}=\tuple{\Gn{!A_1},\dots,\Gn{!A_n}}$.

اگر $\Gamma\Sequent\Delta$ سکوئنت باشد، یک عدد گودلِ
$\Gamma\Sequent\Delta$ عبارت است از
\[
\Gn{\Gamma \Sequent \Delta} = \tuple{\Gn{\Gamma}, \Gn{\Delta}}.
\]

اگر $\pi$ !!a{derivation}ی در~$\Log{LK}$ باشد، $\Gn{\pi}$ چنین تعریف
می‌شود:
\begin{enumerate}
\item اگر $\pi$ فقط از سکوئنت آغازین~$\Gamma\Sequent\Delta$ تشکیل
  شده باشد، آنگاه $\Gn{\pi}$ برابر است با
  \[
    \tuple{0, \Gn{\Gamma \Sequent \Delta}}.
  \]
\item اگر $\pi$ با استنتاجی دارای یک یا دو مقدمه پایان یابد،
  $\Gamma\Sequent\Delta$ نتیجهٔ آن باشد، و~$\pi_1$ و~$\pi_2$
  زیراثبات‌های بی‌واسطه‌ای باشند که به مقدمات آخرین استنتاج منتهی
  می‌شوند، آنگاه $\Gn{\pi}$ به‌ترتیب برابر است با
  \begin{align*}
    & \tuple{1, \Gn{\pi_1}, \Gn{\Gamma \Sequent \Delta}, k} \text{ یا}\\
    & \tuple{2, \Gn{\pi_1}, \Gn{\pi_2}, \Gn{\Gamma \Sequent \Delta}, k},
  \end{align*}
  که در آن~$k$، بر حسب قاعدهٔ به‌کاررفته در آخرین استنتاج، از جدول
  زیر به دست می‌آید:

  \begin{tabular}{lccccccc}
    \text{قاعده:} & \LeftR{\Weakening} & \RightR{\Weakening} &
    \LeftR{\Contraction} & \RightR{\Contraction} &
    \LeftR{\Exchange} & \RightR{\Exchange} \\
    $k$: & 1 & 2 & 3 & 4 & 5 & 6 \\[2ex]
    \text{قاعده:} & \LeftR{\lnot} & \RightR{\lnot} &
    \LeftR{\land} & \RightR{\land} &
    \LeftR{\lor} & \RightR{\lor} \\
    $k$: & 7 & 8 & 9 & 10 & 11 & 12 \\[2ex]
    \text{قاعده:} & \LeftR{\lif} & \RightR{\lif} &
    \LeftR{\lforall} & \RightR{\lforall} &
    \LeftR{\lexists} & \RightR{\lexists} \\
    $k$: & 13 & 14 & 15 & 16 & 17 & 18 \\[2ex]
    \text{قاعده:} & \Cut & = \\
    $k$: & 19 & 20
  \end{tabular}
\end{enumerate}
\end{defn}

\begin{ex}
  !!{derivation} بسیار سادهٔ زیر را در نظر بگیرید:
  \begin{prooftree}
    \Axiom$!A \fCenter !A$
    \RightLabel{\LeftR{\land}}
    \UnaryInf$!A \land !B \fCenter !A$
    \RightLabel{\RightR{\lif}}
    \UnaryInf$\fCenter (!A \land !B) \lif !A$
  \end{prooftree}
  عدد گودل سکوئنت آغازین برابر
  $p_0=\tuple{0,\Gn{!A\Sequent !A}}$ خواهد بود. عدد گودلِ
  !!{derivation} منتهی به نتیجهٔ~$\LeftR{\land}$ برابر
  $p_1=\tuple{1,p_0,\Gn{!A\land !B\Sequent !A},9}$ خواهد بود
  ($1$ چون~$\LeftR{\land}$ یک مقدمه دارد، سپس عدد گودل نتیجهٔ
  $!A\land !B\Sequent !A$، و~$9$ عدد رمزکنندهٔ~$\LeftR{\land}$ است).
  بنابراین عدد گودل کل !!{derivation} برابر
  $\tuple{1,p_1,\Gn{\Sequent(!A\land !B)\lif !A},14}$ است؛ یعنی
  \[
  \tuple{1, \tuple{1, \tuple{0, \Gn{!A \Sequent !A}}, \Gn{!A \land !B \Sequent !A}, 9},
    \Gn{\Sequent (!A \land !B) \lif !A}, 14}.
  \]
\end{ex}

\begin{explain}
پس از تثبیت بازنمایی !!{derivation}s، باید نشان دهیم می‌توانیم چنین
!!{derivation}sیی را به‌طور بازگشتی اولیه دست‌کاری کنیم و ویژگی‌ها و
روابط اساسی‌شان را نیز چنین بیان کنیم. بعضی عملیات ساده‌اند: مثلاً
با داشتن عدد گودل~$p$ از !!a{derivation}ی،
$\fn{EndSeq}(p)=(p)_{(p)_0+1}$ عدد گودل سکوئنت پایانی و
$\fn{LastRule}(p)=(p)_{(p)_0+2}$ رمز آخرین قاعده را به دست می‌دهد.
ویژگی~$\fn{Sequent}(s)$ که با
\[
  \len{s} = 2 \land \bforall{i<\len{(s)_0} + \len{(s)_1}}{\fn{Sent}(((s)_0 \concat (s)_1)_i)}
\]
تعریف می‌شود، دربارهٔ~$s$ برقرار است اگر و تنها اگر $s$ عدد گودل
سکوئنتی متشکل از !!{sentence}s باشد. بعضی موارد بسیار دشوارترند. دست‌کم
طرحی از چگونگی انجامشان می‌دهیم. هدف، نشان‌دادن این است که رابطهٔ
«$\pi$ !!a{derivation}ی از~$!A$ بر پایهٔ~$\Gamma$ است» رابطه‌ای بازگشتی
اولیه میان اعداد گودل~$\pi$ و~$!A$ است.
\end{explain}

\begin{prop}\ollabel{prop:followsby}
  ویژگی~$\fn{Correct}(p)$ که برقرار است اگر و تنها اگر آخرین استنتاج
  در !!{derivation}~$\pi$ با عدد گودل~$p$ درست باشد، بازگشتی اولیه است.
\end{prop}

\begin{proof}
  $\Gamma\Sequent\Delta$ سکوئنت آغازین است اگر یا !!a{sentence}ای چون
  $!A$ وجود داشته باشد که~$\Gamma\Sequent\Delta$ همان
  $!A\Sequent !A$ باشد، یا ترمی چون~$t$ وجود داشته باشد که
  $\Gamma\Sequent\Delta$ همان~$\emptyset\Sequent\eq[t][t]$ باشد. بر
  حسب اعداد گودل، $\fn{InitSeq}(s)$ برقرار است اگر و تنها اگر
  \begin{align*}
  \bexists{x < s}{} (\fn{Sent}(x) & \land
  s = \tuple{\tuple{x},\tuple{x}})
  \lor {}\\
  \bexists{t<s}{} (\fn{Term}(t) & \land
  s = \tuple{0, \tuple{\Gn{{\eq}(} \concat t \concat \Gn{,} \concat t \concat \Gn{)}}}).
  \end{align*}

  همچنین باید نشان دهیم به ازای هر قاعدهٔ استنتاج~$R$، رابطهٔ
  $\fn{FollowsBy}_R(p)$ بازگشتی اولیه است؛
  $\fn{FollowsBy}_R(p)$ برقرار است اگر و
  تنها اگر $p$ عدد گودل !!{derivation}~$\pi$ باشد و سکوئنت پایانی~$\pi$
  با کاربرد درست~$R$ از زیر!!{derivation}sی بی‌واسطهٔ~$\pi$ نتیجه شود.

  حالت ساده، قاعدهٔ~$\RightR{\land}$ است. اگر $\pi$ با استنتاج درست
  $\RightR{\land}$ پایان یابد، چنین شکلی دارد:
  \begin{prooftree}
    \AxiomC{}
    \RightLabel{$\pi_1$}
    \Deduce$\Gamma \fCenter \Delta, !A$

    \AxiomC{}
    \RightLabel{$\pi_2$}
    \Deduce$\Gamma \fCenter \Delta, !B$

    \RightLabel{\RightR\land}
    \BinaryInf$\Gamma \fCenter \Delta, !A \land !B$
  \end{prooftree}
  پس آخرین استنتاج در !!{derivation}~$\pi$ کاربردی درست از~$\RightR{\land}$ است اگر
  و تنها اگر دنباله‌های !!{sentence}sیِ~$\Gamma$ و~$\Delta$ و دو مورد از
  !!{sentence}s، یعنی~$!A$ و~$!B$ وجود داشته باشند به‌گونه‌ای که سکوئنت پایانی
  $\pi_1$ برابر~$\Gamma\Sequent\Delta,!A$، سکوئنت پایانی~$\pi_2$
  برابر~$\Gamma\Sequent\Delta,!B$، و سکوئنت پایانی~$\pi$ برابر
  $\Gamma\Sequent\Delta,!A\land !B$ باشد. فقط باید این را به اعداد
  گودل ترجمه کنیم. اگر $s=\Gn{\Gamma\Sequent\Delta}$، آنگاه
  $(s)_0=\Gn{\Gamma}$ و~$(s)_1=\Gn{\Delta}$. بنابراین
  $\fn{FollowsBy}_{\RightR{\land}}(p)$ برقرار است اگر و تنها اگر
\begin{align*}
& \bexists{g < p}{\bexists{d < p}{\bexists{a < p}{\bexists{b < p}{\quad}}}} \\
& \qquad \fn{EndSeq}(p) =
  \tuple{g, d \concat
    \tuple{\Gn{(} \concat a \concat \Gn{\land} \concat b \concat \Gn{)}}} \land {} \\
& \qquad \fn{EndSeq}((p)_1) = \tuple{g, d \concat \tuple{a}} \land {}\\
& \qquad \fn{EndSeq}((p)_2) = \tuple{g, d \concat \tuple{b}} \land {}\\
& \qquad (p)_0 = 2 \land \fn{LastRule}(p) = 10.
\end{align*}
سطرها به‌ترتیب بیان می‌کنند «دنباله‌ای~($\Gamma$) با عدد گودل~$g$،
دنباله‌ای~($\Delta$) با عدد گودل~$d$، !!a{formula}ی~($!A$) با عدد
گودل~$a$ و !!a{formula}ی~($!B$) با عدد گودل~$b$ وجود دارند» به‌گونه‌ای
که «سکوئنت پایانی~$\pi$ برابر~$\Gamma\Sequent\Delta,!A\land !B$»،
«سکوئنت پایانی~$\pi_1$ برابر~$\Gamma\Sequent\Delta,!A$»، «سکوئنت
پایانی~$\pi_2$ برابر~$\Gamma\Sequent\Delta,!B$»، و «$\pi$ دو
زیر‌اشتقاق بی‌واسطه دارد و آخرین قاعدهٔ استنتاج
$\RightR\land$ (با شمارهٔ~$10$) است».

آخرین استنتاج در~$\pi$ کاربردی درست از~$\RightR\lexists$ است اگر و
تنها اگر دنباله‌های~$\Gamma$ و~$\Delta$، !!a{formula}~$!A$، متغیر~$x$
و ترم بستهٔ~$t$ وجود داشته باشند به‌گونه‌ای که سکوئنت پایانی~$\pi$ برابر
$\Gamma\Sequent\Delta,\lexists[x][!A]$ و سکوئنت پایانی~$\pi_1$
برابر~$\Gamma\Sequent\Delta,\Subst{!A}{t}{x}$ باشد. پس بر حسب اعداد
گودل، $\fn{FollowsBy}_{\RightR\lexists}(p)$ برقرار است اگر و تنها اگر
\begin{align*}
  & \bexists{g<p}{\bexists{d<p}{\bexists{a<p}{\bexists{x<p}{\bexists{t<p}{\quad}}}}}\\
  & \qquad \fn{Var}(x) \land \fn{ClTerm}(t) \land {}\\
  & \qquad \fn{EndSeq}(p) =
  \tuple{
    g,
    d \concat \tuple{\Gn{\lexists} \concat x \concat a}
  } \land {}\\
  & \qquad \fn{EndSeq}((p)_1) =
  \tuple{
    g,
    d \concat \tuple{\fn{Subst}(a, t, x)}
  } \land {}\\
  & \qquad (p)_0 = 1 \land \fn{LastRule}(p) = 18.
\end{align*}
سپس $\fn{Correct}(p)$ را چنین تعریف می‌کنیم:
\begin{multline*}
  \fn{Sequent}(\fn{EndSeq}(p)) \land {}\\
  [(\fn{LastRule}(p) = 1 \land
  \fn{FollowsBy}_{\LeftR\Weakening}(p)) \lor \dots \lor {}\\
  (\fn{LastRule}(p) = 20 \land \fn{FollowsBy}_{\eq}(p)) \lor {}\\
  (p)_0 = 0 \land \fn{InitSeq}(\fn{EndSeq}(p))]
\end{multline*}
سطر نخست تضمین می‌کند سکوئنت پایانی~$p$ واقعاً سکوئنتی متشکل از
!!{sentence}sست. سطر آخر حالتی را پوشش می‌دهد که $p$ فقط یک سکوئنت
آغازین است.
\end{proof}

\begin{prob}
  ویژگی‌های زیر را مانند \olref[inc][art][plk]{prop:followsby} تعریف
  کنید:
  \begin{enumerate}
  \item $\fn{FollowsBy}_{\Cut}(p)$؛
  \item $\fn{FollowsBy}_{\LeftR\lif}(p)$؛
  \item $\fn{FollowsBy}_{\eq}(p)$؛
  \item $\fn{FollowsBy}_{\RightR{\lforall}}(p)$.
  \end{enumerate}
  برای مورد آخر باید همچنین نشان دهید می‌توان به‌طور بازگشتی اولیه
  آزمود که آخرین استنتاجِ !!{derivation} با عدد گودل~$p$ شرط متغیر ویژه
  را برآورده می‌کند؛ یعنی متغیر ویژهٔ~$a$ برای~$\RightR{\lforall}$
  در سکوئنت پایانی رخ نمی‌دهد.
  \end{prob}


\begin{prop}
  \ollabel{prop:deriv}
  رابطهٔ~$\fn{Deriv}(p)$ که برقرار است اگر $p$ عدد گودل یک
  !!{derivation} درست~$\pi$ باشد، بازگشتی اولیه است.
\end{prop}

\begin{proof}
  !!^a{derivation}~$\pi$ درست است اگر هر استنتاج آن کاربردی درست از یک
  قاعده باشد؛ یعنی اگر هر زیر!!{derivation}s آن با استنتاجی درست پایان
  یابد. پس $\fn{Deriv}(p)$ اگر و تنها اگر
  \[
  \bforall{i<\len{\fn{SubtreeSeq}(p)}}{\fn{Correct}((\fn{SubtreeSeq}(p))_i)}.
  \]
\end{proof}


\begin{prop}
فرض کنید $\Gamma$ مجموعه‌ای بازگشتی اولیه از !!{sentence}s باشد. آنگاه
رابطهٔ~$\Prf[\Gamma](x,y)$ که بیان می‌کند «$x$ رمز !!a{derivation}~$\pi$
از~$\Gamma_0\Sequent !A$ برای بعضی~$\Gamma_0\subseteq\Gamma$ متناهی
است و~$y$ عدد گودل~$!A$ است»، بازگشتی اولیه است.
\end{prop}

\begin{proof}
فرض کنید «$y\in\Gamma$» با محمول بازگشتی اولیهٔ~$R_\Gamma(y)$ داده
شود. باید نشان دهیم $\Prf[\Gamma](x,y)$، که برقرار است اگر و تنها اگر
$y$ عدد گودل جملهٔ~$!A$ و~$x$ رمز یک !!{derivation}~$\Log{LK}$ با
سکوئنت پایانی~$\Gamma_0\Sequent !A$ باشد، بازگشتی اولیه است.

بنا بر گزارهٔ پیشین، ویژگی~$\fn{Deriv}(x)$ که برقرار است اگر و تنها
اگر $x$ رمز !!{derivation} درست~$\pi$ در~$\Log{LK}$ باشد بازگشتی اولیه
است. اگر $x$ چنین رمزی باشد، $\fn{EndSeq}(x)$ رمز سکوئنت پایانی
$\pi$ است؛ پس $(\fn{EndSeq}(x))_0$ رمز سمت چپ سکوئنت پایانی و
$(\fn{EndSeq}(x))_1$ رمز سمت راست آن است. بنابراین می‌توانیم
«سمت راست سکوئنت پایانی~$\pi$ برابر~$!A$ است» را با
$\len{(\fn{EndSeq}(x))_1}=1\land
((\fn{EndSeq}(x))_1)_0=y$ بیان کنیم. سمت چپ سکوئنت پایانی~$\pi$
البته خودبه‌خود متناهی است؛ فقط باید بیان کنیم هر جملهٔ آن در~$\Gamma$
است. پس می‌توانیم $\Prf[\Gamma](x,y)$ را چنین تعریف کنیم:
\begin{align*}
\Prf[\Gamma](x, y) \defiff {}&
\fn{Deriv}(x) \land {} \\
& \bforall{i <
  \len{(\fn{EndSeq}(x))_0}}{R_\Gamma(((\fn{EndSeq}(x))_0)_i)} \land {}\\
& \len{(\fn{EndSeq}(x))_1} = 1 \land ((\fn{EndSeq}(x))_1)_0 = y.
\end{align*}
\end{proof}

\end{document}
